
\section{Portfolio backtesting}\label{backtesting}

Backtesting, also called historical simulation or (recently borrowed from Machine Learning terminology) out-of-sample testing
is the process of testing and measuring the predictive power of an investment strategy using historical data that are not part of the calibration set.

In our case we have built an application to simulate, using historical data, an investment in a portfolio periodically rebalanced according with any of the
optimal strategies mentioned in this paper. The weights are calibrated using a predefined historical period, prior to each of the rebalancing dates.
The actual implementation can be accessed through \azapy\ package. In this paper, several simple Python script examples of backtesting were presented for Sharpe type of optimal
portfolios. They can be easily extended to all other optimization strategies (see the documentation at \azapydoc).

We had in mind a retail investor with limited computational resources and trading execution capabilities.\footnote{In contrast to an
institutional investor with ultra-fast computation and execution capabilities. In this case most of the financial slippages mentioned here can be eliminated.}
The portfolio is rebalanced quarterly, 3 business days before the calendar quarter end. The trading execution is assumed to
be performed at the closing prices.
The portfolio weights are computed after the close of the business day before the
execution day.
We have used the New York Stock Exchange business calendar.
The offset of rebalancing date, the offset of weights calculation date, the length of the historical calibration period,
the length of the entire historical simulation, etc. are parameters that can be adjusted in the \azapy\ implementation, see the documentation at \azapydoc.


In this process there are several assumptions that in real life may lead to cash slippages. Let us list them.
\begin{itemize}
  \item The dividends are assumed to be available for reinvestment on the next rebalancing day after the ex-dividend date.
  In general, there is a delay of a week or more for the cash dividend to be collected. Therefore, it is possible that not all 
  cash dividends will be available for reinvestment on the rebalancing day. However, they will be available for the
  next rebalancing day. This will create a reinvestment slippage that, we believe, is small enough to be
  irrelevant to the overall investment performance.
  \item The delay of one business day between the weights (and the number of shares) calculation and the trading execution. As we mentioned, we
  are considering a realistic rebalancing schedule for a regular retail investor. Therefore, we assume that the new positions
  (\ie\ the number of shares in the portfolio components) were calculated using the closing prices of the business day before the rebalancing day.
  Therefore, there is a price slippage between the computation and the actual trading execution prices.
  In general, the slippage is small, since often a rebalancing consists in readjusting existing positions by selling and buying
  in roughly equal cash amounts. Although small, we keep track of this type of slippage in our investment model by
  readjusting the capital at each rebalancing date. Its impact to the portfolio performance is small. \azapy\ implementation
  offer a facility to view and keep track of these adjustments.
  \item Execution price slippage between the actual trading execution and the closing prices. The model simulation assumes that all
  the rebalancing trades are occurring at the closing prices. In real life, a typical retail investor may execute trades only through manually entered orders.
  Although these orders can be placed in the last minutes of the trading day (say last 15 min) it cannot be assumed that they are executed
  at the closing price of the day. This type of slippage is small enough to have no impact to the overall investment
  performance.\footnote{Although not yet popular, several brokers offer Market-on-Close (MOC) orders for retail investors,
  further reducing this type of  slippage.}
  \item Rounding off to the nearest whole number of shares slippage. The model assumes that the portfolio positions are whole numbers of shares.
  Although some brokers accept fractional orders, at the time when we are drafting this paper, it is not yet a widely available offer.
  Our investment model keeps track of the resulting roundoff cash slippage as an adjustment to the capital on each rebalancing date.
  In general, its impact to the portfolio performance is small.\footnote{Assuming that the share prices are
  small relative to the total invested capital.} \azapy\ implementation
  offer a facility to view and keep track of these adjustments.
  \item We do not consider the impact of transaction costs. These days, all major brokers are waiving the transaction fees under reasonable
  conditions for most of the US equity products (single shares and ETF's). If you are using a broker that is charging
  transactions fees or trading in exotic products with transaction fees (\eg\ leveraged ETF's, etc.)
  or using a margin account with borrowing costs, etc., then you should consider the impact of
  these fees.
\end{itemize}

\BackToContents
